\section{Further Work \& Applications}
\label{sec:further-work}

\subsection{Implicit Polymorphism}
\label{sec:future-implicit}

Globally inferred languages typically have much more confusing typing derivations and type errors than locally inferred languages. Extending the calculus to implicit polymorphism in a bidirectional setting would expand the application significantly. In particular, the ability to explain even well-typed code is particularly useful here as error marks are often localised to the wrong location, resulting the real error being in well-typed code.

\subsection{Applications}
\label{sec:applications}

\subsubsection{LLM-Guided Repair}
\label{sec:app-llm}

Contemporary editor-integrated LLM workflows typically view the whole program or a arbitrarily chosen subset and may rewrite unrelated code when asked to fix a single type error. Hole-based LLM tooling such as ChatLSP~\cite{ChatLSP} already support treating incomplete programs directly within an LLM. Type slicing can minimise the incomplete program that the model works with to restrict edits to subterms that provably may affect a queried type error. This applies beyond just type errors, consider a product type for tables, then the LLM could use type slicing to minimally constrain its edits to a particular field of the table. However, remember that this constraint does not imply that other types not present in the query will be unaffected due to non-exactness of minimal slices in the general case (\cref{lem:exact-not-exist}).

For this particular application, edits to change the queried sub-type of a program would not be possible merely with minimal type slices, as these edits must adapt both branches of case statements. The \emph{contribution slice} (\cref{sec:join-slices}) are the more applicable tool here, which are a join of all minimal slices (but can be calculated directly, rather than via minimal slices).

Confining the LLM to the contribution slice gives a verified, minimally-sufficient editing context. The model can query the parts of the type it wishes to change. 

Similarly, when an LLM is understanding a type error, being able to use the type slicing tools in the same way a human would would provide more concise program slices to interpret from, saving tokens and potentially avoid confusion by omitting superfluous information.

\subsubsection{Cast Slicing \& Dynamics}
\label{sec:cast-dynamics}

A gradually-typed language elaborates source programs into a cast calculus and runs the resulting casts dynamically. Type slices can be attached directly to those casts. Hence, casts can then explain why they were inserted by pointing back to code; the source will be a slice of the term being cast, and the target a slice of its context. Approximate minimal slices can be tagged directly onto types allowing them to be decomposed at runtime, allowing dynamic casts to point back to code. This is particularly useful for debugging dynamic type errors, which do not typically point back to code, but do when carrying type slices. However, the strong properties of slices being slices of the original program is not maintained during dynamic execution. Tracking dynamic execution, perhaps integrating ideas from Perera~\cite{Perera2012,PereraGarg2016}, would allow users to trace back casts to source code with stronger guarantees. This application is detailed thoroughly in Carroll~\cite{Carroll2024}.

\subsubsection{Smaller Error Marks}
\label{sec:app-largest-valid}

An error mark in the style of \cref{sec:marking}~\cite{Marking2024} replaces entire inconsistent expressions with dynamic-typed holes, discarding any locally-good type information. For example, when the type inconsistency lies deep in the syntax tree of a type (consider both branches of an if statement being products, but with different projections). Then treating the whole expression as a hole will result in fewer cast insertions than possible making execution more dynamic than needed, resulting in dynamic errors to be caught later (or missed entirely).

\section{Conclusions}
\label{sec:conclusions}

This dissertation extends typing slicing theory to explicit System~F polymorphism, products, and sums; mechanises the statics, the lattice structure, derivation and synthesis slices, the two static gradual guarantee theorems and synthesis unicity, in Agda, reformulating much of the prior theory. A more rigorous account of \emph{assumption slices} was considered and extensive exploration of provably correct algorithms, rather than previous approximations, including various optimisations.

The theory is formally integrated with the error-marking calculus of \textcite{Marking2024}, formally connecting it to ill-typed programs, justifying its \textit{sound} use in debugging \textit{all} static type errors of a program. Finally, type slicing was related to other slicing methods, and further applications were proposed \cref{sec:further-work}.

%%% Local Variables:
%%% mode: LaTeX
%%% TeX-master: "PolymorphicTypeSlicing"
%%% End:
